\documentclass[11pt,a4paper]{article}

\usepackage[T1]{fontenc}
\usepackage[utf8]{inputenc}
\usepackage[brazil]{babel}
\usepackage{geometry}
\geometry{margin=2.5cm,headheight=14pt}
\usepackage{xcolor}
\usepackage{listings}
\usepackage{hyperref}
\usepackage{enumitem}
\usepackage{tikz}
\usetikzlibrary{arrows.meta, positioning}
\usepackage{titlesec}
\usepackage{fancyhdr}

\hypersetup{
    colorlinks=true,
    linkcolor=blue!50!black,
    urlcolor=blue!50!black,
    citecolor=blue!50!black,
    pdftitle={Revisão de Android: Compose, ContentProvider, SharedPreferences, Lifecycle, Coroutines, MVVM, StateFlow/LiveData},
    pdfauthor={Pedro Oliveira}
}

\definecolor{codebg}{RGB}{245,245,245}
\definecolor{codekeyword}{RGB}{127,0,85}
\definecolor{codestring}{RGB}{0,128,0}
\definecolor{codecomment}{RGB}{128,128,128}
\definecolor{codeannotation}{RGB}{102,51,153}

\lstdefinelanguage{Kotlin}{
    comment=[l]{//},
    morecomment=[s]{/*}{*/},
    string=[b]",
    morestring=[b]',
    morekeywords={
        package, import, class, object, interface, fun, val, var, by,
        if, else, when, for, while, do, return, private, public,
        internal, override, companion, data, sealed, enum, in, is,
        null, true, false, this, super, it, modifier
    },
}

\lstdefinestyle{kotlinstyle}{
    language=Kotlin,
    backgroundcolor=\color{codebg},
    basicstyle=\ttfamily\footnotesize,
    keywordstyle=\color{codekeyword}\bfseries,
    commentstyle=\color{codecomment}\itshape,
    stringstyle=\color{codestring},
    numbers=left,
    numberstyle=\tiny\color{gray},
    numbersep=8pt,
    frame=single,
    rulecolor=\color{gray!40},
    breaklines=true,
    breakatwhitespace=false,
    showstringspaces=false,
    tabsize=4,
    columns=flexible,
    captionpos=b,
    literate=
        {ã}{{\~a}}1 {õ}{{\~o}}1 {ç}{{\c c}}1 {á}{{\'a}}1 {é}{{\'e}}1
        {í}{{\'i}}1 {ó}{{\'o}}1 {ú}{{\'u}}1 {â}{{\^a}}1 {ê}{{\^e}}1
        {ô}{{\^o}}1
}

\lstset{style=kotlinstyle}
\emergencystretch=3em
\sloppy

\pagestyle{fancy}
\fancyhf{}
\fancyhead[L]{Revisão de Android — ComposeReview}
\fancyhead[R]{\thepage}
\renewcommand{\headrulewidth}{0.4pt}

\title{\textbf{Revisão de Android} \\ \large Compose, ContentProvider, SharedPreferences, Lifecycle, \\ Coroutines, MVVM e StateFlow/LiveData \\[4pt] \normalsize Aplicação de exemplo: \texttt{ComposeReview}}
\author{Pedro Oliveira}
\date{\today}

\begin{document}
\shorthandoff{"}

\maketitle
\tableofcontents
\newpage

\section{Introdução}

Este documento acompanha o projeto Android \texttt{jetpack-compose-review/},
uma aplicação de exemplo que revisa, na prática, sete tópicos que costumam
cair em entrevista/prova de Android:

\begin{itemize}
    \item Jetpack Compose: \texttt{@Composable}, composição e
    recomposição, \texttt{state}, \texttt{remember},
    \texttt{mutableStateOf} e UI declarativa;
    \item Content Provider e compartilhamento de dados;
    \item SharedPreferences e suas limitações;
    \item Activity/Fragment lifecycle;
    \item Coroutines;
    \item MVVM;
    \item StateFlow / LiveData.
\end{itemize}

O projeto está dividido em duas partes. A \textbf{Parte I}
(Seção~\ref{sec:compose}) cobre só Jetpack Compose puro, com um contador e
uma lista de tarefas cujo estado vive inteiramente na função
\texttt{@Composable} — a versão mais simples possível dos conceitos de
Compose. A \textbf{Parte II} (Seções~\ref{sec:contentprovider} a
\ref{sec:stateflow-livedata}) adiciona uma terceira tela, \emph{Notas},
que já não guarda estado só localmente: ela persiste dados em SQLite por
trás de um \texttt{ContentProvider}, guarda preferências simples em
\texttt{SharedPreferences}, organiza a lógica em camadas MVVM, usa
\texttt{Coroutines} para I/O assíncrono e expõe o estado da tela via
\texttt{StateFlow} (comparado lado a lado com \texttt{LiveData}). Uma
quarta tela, separada e escrita com Views/XML clássicos, existe só para
tornar visível o \texttt{Activity}/\texttt{Fragment lifecycle}
completo, algo que não aparece em um app 100\% Compose de Activity única.

\section{Estrutura do projeto}

O projeto segue a estrutura padrão de um módulo Android/Gradle:

\begin{verbatim}
jetpack-compose-review/
|-- build.gradle.kts                  # configuracao raiz (plugins)
|-- settings.gradle.kts               # declara o modulo :app
|-- app/
|   |-- build.gradle.kts              # dependencias (Compose, lifecycle,
|   |                                 #   coroutines, appcompat/fragment)
|   `-- src/main/
|       |-- AndroidManifest.xml       # registra Activities e o Provider
|       |-- java/com/pedro/composereview/
|       |   |-- MainActivity.kt       # Activity Compose + telas 1 e 2
|       |   |-- data/                 # Parte II: ContentProvider
|       |   |   |-- Note.kt
|       |   |   |-- NotesContract.kt
|       |   |   |-- NotesDbHelper.kt      (SQLiteOpenHelper)
|       |   |   |-- NotesProvider.kt      (ContentProvider)
|       |   |   `-- NotesRepository.kt    (Coroutines + Flow)
|       |   |-- prefs/
|       |   |   `-- UserPreferences.kt    (SharedPreferences)
|       |   |-- ui/
|       |   |   |-- notes/
|       |   |   |   |-- NotesViewModel.kt (MVVM + StateFlow/LiveData)
|       |   |   |   `-- NotesScreen.kt    (tela 3, Compose)
|       |   |   `-- theme/
|       |   |       |-- Color.kt
|       |   |       |-- Theme.kt
|       |   |       `-- Type.kt
|       |   `-- lifecycle/            # Parte II: demo classica
|       |       |-- LifecycleDemoActivity.kt
|       |       `-- LifecycleDemoFragment.kt
|       `-- res/
|           |-- layout/               # XML só das telas de lifecycle
|           |   |-- activity_lifecycle_demo.xml
|           |   `-- fragment_lifecycle_demo.xml
|           `-- values/
|               |-- strings.xml
|               `-- themes.xml
`-- docs/
    `-- revisao-jetpack-compose.tex   # este documento
\end{verbatim}

A regra geral do projeto: tudo que é tela normal do app é Compose puro
(\texttt{MainActivity.kt} e o pacote \texttt{ui/}); \texttt{res/layout/} e
XML existem \emph{apenas} na demonstração de lifecycle clássico, isolada
no pacote \texttt{lifecycle/}, para não misturar os dois estilos na
mesma tela.

\section{Parte I --- Jetpack Compose}
\label{sec:compose}

\subsection{Os seis conceitos, um a um}

\subsubsection{\texttt{@Composable}}

\texttt{@Composable} é a anotação que transforma uma função Kotlin comum
em uma \textbf{unidade de UI} que o compilador do Compose sabe como
rastrear, (re)executar e otimizar. Uma função marcada com
\texttt{@Composable} pode chamar outras funções \texttt{@Composable}, ler
estado (\texttt{State<T>}) e emitir elementos de UI — mas não pode ser
chamada de código Kotlin comum fora desse contexto.

No projeto, todas as telas são funções \texttt{@Composable}:
\texttt{ReviewApp}, \texttt{SectionTitle}, \texttt{CounterCard},
\texttt{TodoCard}, \texttt{TaskList} e \texttt{TaskRow}. Exemplo:

\begin{lstlisting}[caption={Uma função @Composable simples e reutilizável}, label=lst:composable]
@Composable
private fun SectionTitle(text: String) {
    Text(
        text = text,
        fontWeight = FontWeight.Bold,
        style = MaterialTheme.typography.titleMedium
    )
}
\end{lstlisting}

Note que \texttt{SectionTitle} não devolve uma \texttt{View}: ela apenas
\emph{descreve} que, naquele ponto da árvore, deve existir um
\texttt{Text} com determinado estilo. Quem decide \emph{quando} e
\emph{como} desenhar isso é o runtime do Compose — esse é o cerne da UI
declarativa (Seção~\ref{sec:declarativa}).

\subsubsection{Composição e recomposição}

\textbf{Composição} é o processo de executar as funções
\texttt{@Composable} e montar a árvore de UI resultante (quais elementos
existem, com quais parâmetros, aninhados de que forma). Isso acontece na
primeira renderização da tela.

\textbf{Recomposição} é o processo de executar \emph{de novo} apenas as
funções \texttt{@Composable} cujas leituras de estado mudaram, para
atualizar a árvore — sem recriar a \texttt{Activity} nem reconstruir a
UI do zero. O Compose rastreia, para cada composable, quais objetos de
\texttt{State} ele leu durante a composição; quando um desses valores
muda, apenas os composables que o leram (e, quando necessário, seus
descendentes) são recompostos.

Para tornar esse processo visível, \texttt{CounterCard} mantém um
contador de recomposições (\texttt{recompositions}) que \emph{não} é um
\texttt{State} — é uma variável comum guardada com \texttt{remember}. Ela
é incrementada toda vez que a função roda de novo, mas sua leitura não
"assina" a composição atual (diferente de um \texttt{mutableStateOf}),
então ela não causa recomposições adicionais por conta própria — apenas
reflete quantas vezes o \texttt{count} mudou até agora:

\begin{lstlisting}[caption={Tornando a recomposição visível na tela}, label=lst:recomposicao]
@Composable
fun CounterCard() {
    var count by remember { mutableStateOf(0) }
    val recompositions = remember { intArrayOf(0) }
    recompositions[0] += 1

    Card(/* ... */) {
        Column(/* ... */) {
            Text("Valor: $count")
            Text("Esta secao recompos ${recompositions[0]} vez(es)")
            Row {
                Button(onClick = { count++ }) { Text("+1") }
                Button(onClick = { count-- }) { Text("-1") }
                TextButton(onClick = { count = 0 }) { Text("Resetar") }
            }
        }
    }
}
\end{lstlisting}

A cada toque em "+1", "-1" ou "Resetar": (1) \texttt{count} muda; (2) o
Compose marca \texttt{CounterCard} como inválido, pois foi ele quem leu
\texttt{count} durante a composição; (3) \texttt{CounterCard} é executado
de novo (recomposição); (4) \texttt{recompositions[0]} é incrementado e o
novo valor aparece na tela — evidenciando que a recomposição de fato
ocorreu.

\subsubsection{\texttt{state}}
\label{sec:state}

Em Compose, \textbf{state} (estado) é qualquer valor que, ao mudar,
deveria fazer a UI se atualizar. Ele é representado pela interface
\texttt{State<T>} (e sua variante mutável, \texttt{MutableState<T>}). O
Compose observa leituras de \texttt{State} feitas dentro de um escopo de
composição e usa isso para saber o que recompor quando o valor muda.

O projeto usa estado em dois lugares com naturezas diferentes:

\begin{itemize}
    \item \textbf{Estado local e primitivo}: \texttt{count} (um
    \texttt{Int}) em \texttt{CounterCard} e \texttt{input} (uma
    \texttt{String}) em \texttt{TodoCard}.
    \item \textbf{Estado de coleção}: \texttt{tasks}, uma
    \texttt{SnapshotStateList<String>} criada com
    \texttt{mutableStateListOf(...)}, que notifica a composição quando
    itens são adicionados ou removidos (\texttt{tasks.add(...)},
    \texttt{tasks.remove(...)}), sem precisar reatribuir a variável
    inteira.
\end{itemize}

Também é demonstrado o padrão de \textbf{state hoisting}: em vez de cada
composable filho guardar seu próprio estado, \texttt{TodoCard} concentra
o estado (\texttt{input} e \texttt{tasks}) e o repassa como parâmetros e
callbacks para \texttt{TaskList} e \texttt{TaskRow}, que ficam
\emph{stateless} — apenas recebem dados e informam eventos para cima:

\begin{lstlisting}[caption={State hoisting: pai com estado, filhos sem estado}, label=lst:hoisting]
@Composable
fun TaskList(tasks: List<String>, onRemove: (String) -> Unit) {
    Column {
        tasks.forEach { task ->
            TaskRow(task = task, onRemove = { onRemove(task) })
        }
    }
}

@Composable
fun TaskRow(task: String, onRemove: () -> Unit) {
    Row {
        Text(text = task, modifier = Modifier.weight(1f))
        TextButton(onClick = onRemove) { Text("Remover") }
    }
}
\end{lstlisting}

Essa separação torna \texttt{TaskList} e \texttt{TaskRow} fáceis de
reutilizar, testar e pré-visualizar isoladamente, porque não dependem de
\emph{onde} o estado mora — apenas do que recebem por parâmetro.

\subsubsection{\texttt{remember}}

Funções \texttt{@Composable} podem ser executadas várias vezes
(recompostas). Sem alguma forma de persistir valores entre essas
execuções, qualquer variável local seria recriada do zero a cada
recomposição — um contador sempre voltaria a zero, por exemplo.

\texttt{remember} resolve isso: ele guarda um valor na memória interna da
composição (a \emph{slot table}) e o devolve nas recomposições
seguintes, em vez de recalculá-lo. O valor só é recriado se o composable
sair completamente da árvore (por exemplo, some da tela) e entrar de
novo, ou se as chaves passadas a variações como \texttt{remember(key)}
mudarem.

No projeto, \texttt{remember} aparece em toda inicialização de estado:

\begin{lstlisting}[caption={remember preservando estado entre recomposições}, label=lst:remember]
var count by remember { mutableStateOf(0) }
val recompositions = remember { intArrayOf(0) }
var input by remember { mutableStateOf("") }
val tasks = remember { mutableStateListOf("Revisar @Composable", "Estudar remember") }
\end{lstlisting}

Em todos os casos, a lambda passada para \texttt{remember} só é executada
\textbf{uma vez} (na primeira composição); nas recomposições seguintes,
\texttt{remember} simplesmente devolve o valor já existente.

\subsubsection{\texttt{mutableStateOf}}

\texttt{mutableStateOf(valorInicial)} cria um objeto
\texttt{MutableState<T>} observável: sempre que sua propriedade
\texttt{.value} é reatribuída, o Compose sabe que precisa invalidar (e
depois recompor) qualquer composable que tenha lido esse valor durante a
composição.

Combinado com o delegate de propriedade \texttt{by} (do pacote
\texttt{androidx.compose.runtime}, via \texttt{getValue}/\texttt{setValue}),
é possível tratar o estado quase como uma variável comum:

\begin{lstlisting}[caption={mutableStateOf com o delegate `by`}, label=lst:mutablestateof]
var count by remember { mutableStateOf(0) }
// equivalente, sem o delegate `by`:
// val countState = remember { mutableStateOf(0) }
// countState.value++

count++          // le e escreve em countState.value internamente
\end{lstlisting}

O mesmo padrão é usado para \texttt{input} em \texttt{TodoCard}
(\texttt{var input by remember \{ mutableStateOf("") \}}), atualizado a
cada tecla digitada via \texttt{onValueChange = \{ input = it \}} do
\texttt{OutlinedTextField}. Para a lista de tarefas, o equivalente para
coleções é usado: \texttt{mutableStateListOf(...)}, que devolve uma lista
cujas operações de mutação (\texttt{add}, \texttt{remove}) já disparam
recomposição sozinhas.

\subsubsection{UI declarativa}
\label{sec:declarativa}

Em UI \textbf{imperativa} (o modelo clássico de \texttt{View}s no
Android), o desenvolvedor descreve \emph{passo a passo como transformar}
a UI de um estado para outro: cria uma \texttt{View}, chama
\texttt{findViewById}, depois \texttt{setText(...)},
\texttt{setVisibility(...)}, etc., sempre que algo muda.

Em UI \textbf{declarativa} (o modelo do Compose), o desenvolvedor
descreve apenas \emph{como a UI deve ser para um dado estado atual} — uma
função pura, em essência: \texttt{UI = f(state)}. Quando o estado muda, a
função é executada de novo (recomposição) e o Compose calcula sozinho a
diferença necessária na árvore de UI, sem que o código precise dizer
"esconda este botão" ou "atualize aquele texto" manualmente.

Isso fica evidente em \texttt{CounterCard}: o texto \texttt{"Valor:
\$count"} não é atualizado com um comando imperativo do tipo
\texttt{textView.setText(...)}. Em vez disso, a função inteira é
reexecutada com o novo valor de \texttt{count}, e a nova string
\texttt{"Valor: \$count"} é simplesmente o resultado dessa execução:

\begin{lstlisting}[caption={UI = f(state): o texto e apenas funcao do estado atual}, label=lst:declarativa]
Text("Valor: $count", style = MaterialTheme.typography.headlineSmall)
\end{lstlisting}

Não existe, em nenhum lugar do código, uma chamada explícita para
"atualizar o texto" — o Compose garante isso automaticamente porque
\texttt{Text} leu \texttt{count} durante a composição.

\subsection{Fluxo de dados: do toque na tela até a recomposição}

O diagrama abaixo resume o ciclo que ocorre a cada toque no botão
\texttt{"+1"} do contador:

\begin{center}
\begin{tikzpicture}[
    node distance=1.1cm and 1.6cm,
    box/.style={draw, rounded corners, minimum width=3.6cm, minimum height=1cm, align=center, font=\small},
    arr/.style={-{Latex[length=2mm]}, thick}
]
\node[box] (tap) {Usuário toca em "+1"};
\node[box, below=of tap] (event) {\texttt{onClick} dispara\\\texttt{count++}};
\node[box, below=of event] (state) {\texttt{mutableStateOf}\\atualiza \texttt{.value}};
\node[box, below=of state] (invalidate) {Compose invalida\\\texttt{CounterCard}};
\node[box, below=of invalidate] (recompose) {\texttt{CounterCard} é\\recomposto};
\node[box, below=of recompose] (ui) {Nova UI: novo\\\texttt{count} e\\\texttt{recompositions[0]}};

\draw[arr] (tap) -- (event);
\draw[arr] (event) -- (state);
\draw[arr] (state) -- (invalidate);
\draw[arr] (invalidate) -- (recompose);
\draw[arr] (recompose) -- (ui);
\end{tikzpicture}
\end{center}

O mesmo ciclo se repete para a lista de tarefas, com a diferença de que
o estado (\texttt{tasks}) mora em \texttt{TodoCard} e é apenas
\emph{lido} por \texttt{TaskList}/\texttt{TaskRow} — por isso, ao
adicionar ou remover uma tarefa, é \texttt{TodoCard} (e os composables
que efetivamente leem \texttt{tasks}) que recompõe, não a árvore inteira
do aplicativo.

\section{Parte II --- arquitetura, dados e ciclo de vida}
\label{sec:parte2}

A Parte I usou estado 100\% local: quando \texttt{MainActivity} morre
(app fechado, processo do sistema encerrado por falta de memória), o
contador e as tarefas somem. A Parte II resolve isso com a tela
\emph{Notas}, construída em camadas — exatamente o tipo de arquitetura
usada em apps reais — e introduz os outros cinco tópicos revisados.

\subsection{Content Provider e compartilhamento de dados}
\label{sec:contentprovider}

\texttt{ContentProvider} é um dos quatro componentes de app do Android
(junto com Activity, Service e BroadcastReceiver). Sua função é expor
dados por trás de um contrato de \textbf{URIs}, para que outros
componentes — ou até outros apps, quando \texttt{exported="true"} — leiam
e escrevam sem conhecer o mecanismo de armazenamento real (SQLite, um
arquivo, uma API remota, etc.). É o mesmo princípio usado pelo Android
para os provedores nativos de Contatos, Calendário e Mídia.

No projeto, \texttt{NotesProvider} expõe a tabela SQLite \texttt{notes}
(criada por \texttt{NotesDbHelper}) através de duas URIs — a coleção
inteira e um item específico — resolvidas com um \texttt{UriMatcher}:

\begin{lstlisting}[caption={NotesProvider: query e insert via URI, com notifyChange}, label=lst:contentprovider]
class NotesProvider : ContentProvider() {

    private lateinit var dbHelper: NotesDbHelper

    override fun onCreate(): Boolean {
        dbHelper = NotesDbHelper(context!!)
        return true
    }

    override fun query(
        uri: Uri, projection: Array<out String>?, selection: String?,
        selectionArgs: Array<out String>?, sortOrder: String?
    ): Cursor {
        val db = dbHelper.readableDatabase
        val cursor = when (uriMatcher.match(uri)) {
            NOTES -> db.query(TABLE_NAME, projection, selection, selectionArgs,
                null, null, sortOrder ?: "created_at DESC")
            NOTE_ID -> db.query(TABLE_NAME, projection, "_id = ?",
                arrayOf(uri.lastPathSegment), null, null, null)
            else -> throw IllegalArgumentException("URI desconhecida: $uri")
        }
        cursor.setNotificationUri(context!!.contentResolver, uri)
        return cursor
    }

    override fun insert(uri: Uri, values: ContentValues?): Uri {
        val db = dbHelper.writableDatabase
        val id = db.insert(TABLE_NAME, null, values)
        val itemUri = ContentUris.withAppendedId(NotesContract.CONTENT_URI, id)
        context!!.contentResolver.notifyChange(itemUri, null)
        return itemUri
    }

    // update(), delete() e getType() seguem o mesmo padrao.
}
\end{lstlisting}

Dois detalhes merecem destaque:

\begin{itemize}
    \item \texttt{cursor.setNotificationUri(...)} e
    \texttt{contentResolver.notifyChange(...)} são o que permite que
    \emph{qualquer} componente que tenha registrado um
    \texttt{ContentObserver} naquela URI seja avisado quando os dados
    mudam — inclusive se a escrita veio de outro processo. É essa
    notificação que \texttt{NotesRepository} usa para transformar
    mudanças no banco em um \texttt{Flow} reativo (Seção~\ref{sec:coroutines}).
    \item \texttt{NotesProvider} está declarado no
    \texttt{AndroidManifest.xml} com \texttt{exported="false"}: só este
    app pode acessá-lo. Compartilhar de fato com outros apps seria
    trocar para \texttt{exported="true"} e, tipicamente, adicionar
    \texttt{readPermission}/\texttt{writePermission} para controlar quem
    pode ler e quem pode escrever.
\end{itemize}

Na prática, o resto do app nunca importa \texttt{NotesDbHelper}
diretamente — só fala com \texttt{ContentResolver} usando
\texttt{NotesContract.CONTENT\_URI}, o que deixa a implementação de
armazenamento livre para mudar sem afetar quem consome os dados.

\subsection{SharedPreferences e suas limitações}
\label{sec:sharedpreferences}

\texttt{SharedPreferences} guarda pares chave/valor simples em um
arquivo XML privado do app — ideal para preferências pequenas (um nome de
usuário, um contador, uma flag de "já mostrou o tutorial"). No projeto,
\texttt{UserPreferences} envolve esse acesso:

\begin{lstlisting}[caption={UserPreferences: leitura/escrita simples com apply()}, label=lst:sharedprefs]
class UserPreferences(context: Context) {
    private val prefs: SharedPreferences =
        context.getSharedPreferences(PREFS_NAME, Context.MODE_PRIVATE)

    fun getUserName(): String = prefs.getString(KEY_USER_NAME, "") ?: ""

    fun setUserName(name: String) {
        prefs.edit().putString(KEY_USER_NAME, name).apply()
    }

    fun incrementNotesCreatedTotal() {
        prefs.edit().putInt(KEY_NOTES_CREATED, getNotesCreatedTotal() + 1).apply()
    }
}
\end{lstlisting}

No app, \texttt{NotesViewModel} usa essa classe para lembrar o nome
digitado pelo usuário e o total histórico de notas criadas — valores que
sobrevivem mesmo que o processo do app seja encerrado, ao contrário do
\texttt{remember} da Parte I.

\texttt{apply()} grava em disco de forma assíncrona (não bloqueia quem
chamou, mas também não confirma que a escrita já terminou);
\texttt{commit()} é a alternativa síncrona, que bloqueia a thread
chamadora até gravar e devolve um \texttt{Boolean} de sucesso — por isso
\texttt{commit()} nunca deveria ser chamado na thread principal.

\textbf{Limitações} que tornam SharedPreferences inadequado para além de
configurações simples:

\begin{itemize}
    \item Chaves são \texttt{String}s soltas — sem verificação do
    compilador, é fácil digitar errado (\texttt{"user\_name"} vs
    \texttt{"username"}) sem nenhum aviso.
    \item Não existe consulta, filtro, relação entre chaves ou schema —
    é só um dicionário chave/valor, não um banco de dados; dados
    estruturados ou relacionais pertencem ao SQLite/ContentProvider
    (Seção~\ref{sec:contentprovider}), não aqui.
    \item O arquivo XML inteiro é carregado para a memória na primeira
    leitura; arquivos grandes custam I/O logo na inicialização do app.
    \item Não é observável de forma nativa por Compose/Flow — só existe
    um \texttt{OnSharedPreferenceChangeListener} manual baseado em
    callback, então não serve como fonte de verdade reativa (comparar
    com \texttt{StateFlow}, Seção~\ref{sec:stateflow-livedata}).
    \item Sem suporte a migração de schema — mudar o formato de um valor
    salvo é responsabilidade inteiramente do código do app.
\end{itemize}

Para os casos em que essas limitações pesam — dados maiores, necessidade
de observar mudanças de forma reativa, ou tipagem mais forte — o Android
recomenda hoje o \textbf{Jetpack DataStore} (Preferences ou Proto), que
expõe os valores como \texttt{Flow} em vez de leitura síncrona.

\subsection{Arquitetura MVVM}
\label{sec:mvvm}

MVVM (\emph{Model-View-ViewModel}) separa a tela em três papéis:

\begin{itemize}
    \item \textbf{Model}: os dados e as regras de acesso a eles —
    aqui, \texttt{Note}, \texttt{NotesRepository} e
    \texttt{UserPreferences}.
    \item \textbf{ViewModel}: dona do estado de UI e da lógica de
    "o que fazer quando o usuário faz X" — aqui,
    \texttt{NotesViewModel}. Sobrevive a mudanças de configuração
    (rotação de tela) porque seu ciclo de vida é gerenciado pelo
    Android Jetpack, não pela Activity/Composable.
    \item \textbf{View}: só renderiza o estado atual e repassa eventos
    do usuário para o ViewModel — aqui, o composable
    \texttt{NotesCard} (Seção~\ref{sec:stateflow-livedata}).
\end{itemize}

\begin{lstlisting}[caption={NotesViewModel: estado + eventos, sem nenhuma referência a Compose}, label=lst:viewmodel]
class NotesViewModel(application: Application) : AndroidViewModel(application) {

    private val repository = NotesRepository(application.contentResolver)
    private val preferences = UserPreferences(application)
    private val _uiState = MutableStateFlow(NotesUiState())
    val uiState: StateFlow<NotesUiState> = _uiState.asStateFlow()

    init {
        viewModelScope.launch {
            repository.observeNotes().collect { notes ->
                _uiState.update { it.copy(notes = notes, isLoading = false) }
            }
        }
    }

    fun addNote(title: String) {
        if (title.isBlank()) return
        viewModelScope.launch {
            repository.addNote(title.trim())
            preferences.incrementNotesCreatedTotal()
        }
    }

    fun deleteNote(id: Long) {
        viewModelScope.launch { repository.deleteNote(id) }
    }
}
\end{lstlisting}

A propriedade central de MVVM aqui é a \textbf{direção única de
dependência}: \texttt{NotesViewModel} não importa nada de
\texttt{androidx.compose}, e \texttt{NotesCard} não sabe nada sobre
\texttt{ContentResolver} ou SQLite — ela só lê \texttt{uiState} e chama
\texttt{addNote}/\texttt{deleteNote}. Isso torna o ViewModel testável
sem precisar renderizar UI, e a UI substituível (poderia virar uma tela
com Views/XML amanhã) sem tocar na lógica de negócio.

\subsection{Coroutines}
\label{sec:coroutines}

Toda operação de I/O no Android (disco, banco, rede) deveria rodar fora
da \emph{main thread}, sob pena de travar a UI (e, em casos extremos,
disparar um ANR — \emph{Application Not Responding}). Coroutines
Kotlin resolvem isso com funções \texttt{suspend} que "pausam" sem
bloquear a thread, mais um \texttt{CoroutineScope} que controla quando
esse trabalho começa e é cancelado.

\texttt{NotesRepository} usa duas técnicas de coroutines:

\begin{lstlisting}[caption={withContext para offload de I/O e callbackFlow para observar mudanças}, label=lst:coroutines]
suspend fun addNote(title: String): Unit = withContext(Dispatchers.IO) {
    val values = ContentValues().apply {
        put(NotesContract.Columns.TITLE, title)
        put(NotesContract.Columns.CREATED_AT, System.currentTimeMillis())
    }
    contentResolver.insert(NotesContract.CONTENT_URI, values)
    Unit
}

fun observeNotes(): Flow<List<Note>> = callbackFlow {
    val observer = object : ContentObserver(Handler(Looper.getMainLooper())) {
        override fun onChange(selfChange: Boolean) { trySend(queryNotes()) }
    }
    contentResolver.registerContentObserver(NotesContract.CONTENT_URI, true, observer)
    trySend(queryNotes())
    awaitClose { contentResolver.unregisterContentObserver(observer) }
}.flowOn(Dispatchers.IO)
\end{lstlisting}

\begin{itemize}
    \item \texttt{withContext(Dispatchers.IO)} move o bloco para uma
    thread pool otimizada para I/O e devolve o resultado quando termina
    — usado em \texttt{addNote}/\texttt{deleteNote}, que fazem uma
    escrita pontual.
    \item \texttt{callbackFlow} faz a ponte entre o mundo de callbacks
    do Android (\texttt{ContentObserver.onChange}) e o mundo de
    coroutines (\texttt{Flow}): cada mudança nos dados chama
    \texttt{trySend(...)}, emitindo uma nova lista para quem estiver
    coletando o \texttt{Flow}. \texttt{awaitClose} garante que o
    \texttt{ContentObserver} é removido quando ninguém mais coleta,
    evitando vazamento.
\end{itemize}

Quem inicia essas coroutines é sempre o \texttt{ViewModel}, com
\texttt{viewModelScope.launch \{ ... \}} (Seção~\ref{sec:mvvm}): esse
escopo é cancelado automaticamente em \texttt{onCleared()}, então uma
coroutine nunca continua rodando depois que a tela que a originou não
existe mais.

\subsection{StateFlow / LiveData}
\label{sec:stateflow-livedata}

Ambos são contêineres \textbf{observáveis} — a UI se inscreve e é
notificada quando o valor muda — mas com origens e comportamento
diferentes:

\begin{itemize}
    \item \textbf{LiveData} é mais antigo (pré-Coroutines/Flow), é
    \emph{lifecycle-aware} por padrão (só entrega valores para
    observadores \texttt{STARTED} ou mais) e é observado com
    \texttt{observe(lifecycleOwner) \{ ... \}} fora do Compose, ou
    \texttt{observeAsState()} dentro dele.
    \item \textbf{StateFlow} é construído sobre Coroutines/Flow, sempre
    tem um valor atual (não existe "ainda não emitiu nada"), suporta os
    operadores de \texttt{Flow} (\texttt{map}, \texttt{combine},
    \texttt{debounce}, ...) e é a opção recomendada atualmente para
    expor estado de um \texttt{ViewModel}.
\end{itemize}

Para comparar os dois lado a lado, \texttt{NotesViewModel} expõe a
mesma informação (quantidade de notas) das duas formas — a LiveData é
derivada do próprio StateFlow com \texttt{asLiveData()}:

\begin{lstlisting}[caption={A mesma informacao exposta como StateFlow e como LiveData}, label=lst:stateflow-livedata]
val uiState: StateFlow<NotesUiState> = _uiState.asStateFlow()

val noteCountLiveData: LiveData<Int> = uiState
    .map { it.notes.size }
    .asLiveData()
\end{lstlisting}

E consumidas na tela (\texttt{NotesCard}) com o equivalente Compose de
cada uma:

\begin{lstlisting}[caption={Coletando StateFlow e observando LiveData a partir do Compose}, label=lst:collect-observe]
val uiState by viewModel.uiState.collectAsStateWithLifecycle()
val liveNoteCount by viewModel.noteCountLiveData.observeAsState(initial = 0)
\end{lstlisting}

\texttt{collectAsStateWithLifecycle()} só mantém a coleta ativa enquanto
a tela está pelo menos \texttt{STARTED}, evitando trabalho desperdiçado
em background — o mesmo comportamento lifecycle-aware que a LiveData já
tinha nativamente.

\subsection{Activity/Fragment lifecycle}
\label{sec:lifecycle}

Uma \texttt{Activity} passa por estados bem definidos —
\texttt{onCreate} $\to$ \texttt{onStart} $\to$ \texttt{onResume} $\to$
\texttt{onPause} $\to$ \texttt{onStop} $\to$ \texttt{onDestroy} — que o
sistema dispara conforme a tela fica visível, ganha foco, perde foco ou
sai de cena. Um \texttt{Fragment} tem um ciclo ainda mais granular,
porque sua \emph{view} pode ser destruída e recriada
(\texttt{onCreateView}/\texttt{onDestroyView}) sem que o Fragment
inteiro seja destruído, e porque ele precisa se anexar/desanexar de uma
Activity hospedeira (\texttt{onAttach}/\texttt{onDetach}).

Como \texttt{MainActivity} é uma Activity única 100\% Compose, esse
ciclo completo de Fragment não apareceria em nenhum outro lugar do
projeto. Por isso existe \texttt{LifecycleDemoActivity} — uma
\texttt{AppCompatActivity} clássica, com layout XML, que hospeda
\texttt{LifecycleDemoFragment} — acessível pelo botão "Abrir demo de
lifecycle" na tela principal:

\begin{lstlisting}[caption={Fragment: cada callback grava e reimprime o log na tela}, label=lst:lifecycle-fragment]
class LifecycleDemoFragment : Fragment(R.layout.fragment_lifecycle_demo) {
    private val entries = mutableListOf<String>()

    override fun onAttach(context: Context) { super.onAttach(context); log("Fragment onAttach") }
    override fun onCreate(savedInstanceState: Bundle?) { super.onCreate(savedInstanceState); log("Fragment onCreate") }
    override fun onViewCreated(view: View, savedInstanceState: Bundle?) {
        super.onViewCreated(view, savedInstanceState); log("Fragment onViewCreated")
    }
    override fun onStart() { super.onStart(); log("Fragment onStart") }
    override fun onResume() { super.onResume(); log("Fragment onResume") }
    override fun onPause() { super.onPause(); log("Fragment onPause") }
    override fun onStop() { super.onStop(); log("Fragment onStop") }
    override fun onDestroyView() { log("Fragment onDestroyView"); super.onDestroyView() }
    override fun onDetach() { super.onDetach(); log("Fragment onDetach") }

    private fun log(message: String) {
        Log.d("LifecycleDemo", message)
        entries += message
        view?.findViewById<TextView>(R.id.text_log)?.text = entries.joinToString("\n")
    }
}
\end{lstlisting}

Rodando a demo, a ordem observada no Logcat (tag \texttt{LifecycleDemo})
é: \texttt{Activity.onCreate} $\to$ \texttt{Fragment.onAttach} $\to$
\texttt{Fragment.onCreate} $\to$ \texttt{Fragment.onViewCreated} $\to$
\texttt{Activity.onStart} $\to$ \texttt{Fragment.onStart} $\to$
\texttt{Activity.onResume} $\to$ \texttt{Fragment.onResume} — deixando
claro que o Fragment está aninhado dentro do ciclo de vida da Activity
que o hospeda, nunca à frente dele.

Mesmo em um app só-Compose, o \texttt{Lifecycle} continua relevante: o
próprio \texttt{NotesCard} observa eventos de \texttt{ON\_START}/
\texttt{ON\_STOP} da Activity para fins de log, usando a ponte oficial
entre Compose e o sistema de Lifecycle clássico:

\begin{lstlisting}[caption={Observando o Lifecycle da Activity a partir de um Composable}, label=lst:lifecycle-compose]
@Composable
private fun ObserveLifecycle(tag: String) {
    val lifecycleOwner = LocalLifecycleOwner.current
    DisposableEffect(lifecycleOwner) {
        val observer = LifecycleEventObserver { _, event ->
            if (event == Lifecycle.Event.ON_START || event == Lifecycle.Event.ON_STOP) {
                Log.d(tag, "Lifecycle event: $event")
            }
        }
        lifecycleOwner.lifecycle.addObserver(observer)
        onDispose { lifecycleOwner.lifecycle.removeObserver(observer) }
    }
}
\end{lstlisting}

\texttt{DisposableEffect} garante que o observador seja removido em
\texttt{onDispose} quando o composable sai da árvore — sem isso, o
observador vazaria e continuaria recebendo eventos de um Lifecycle cuja
tela já não existe mais.

\section{Fluxo de dados da tela Notas: do toque em "Add" até a UI atualizada}

O diagrama abaixo conecta todos os tópicos da Parte II em um único
ciclo — do toque no botão até a lista reaparecer atualizada:

\begin{center}
\begin{tikzpicture}[
    node distance=0.55cm and 1.6cm,
    box/.style={draw, rounded corners, minimum width=8.6cm, minimum height=0.9cm, align=center, font=\small},
    arr/.style={-{Latex[length=2mm]}, thick}
]
\node[box] (tap) {UI (Compose): usuário toca em "Add" $\to$ \texttt{viewModel.addNote(texto)}};
\node[box, below=of tap] (vm) {ViewModel: \texttt{viewModelScope.launch \{ ... \}} (Coroutines)};
\node[box, below=of vm] (repo) {Repository: \texttt{withContext(Dispatchers.IO)} $\to$ \texttt{contentResolver.insert(...)}};
\node[box, below=of repo] (provider) {ContentProvider: \texttt{insert()} grava no SQLite e chama \texttt{notifyChange()}};
\node[box, below=of provider] (observer) {ContentObserver.onChange $\to$ \texttt{callbackFlow} reemite a lista};
\node[box, below=of observer] (state) {ViewModel: \texttt{\_uiState.update \{ ... \}} (StateFlow)};
\node[box, below=of state] (ui) {UI: \texttt{collectAsStateWithLifecycle()} recompõe \texttt{NotesCard} com a nova lista};

\draw[arr] (tap) -- (vm);
\draw[arr] (vm) -- (repo);
\draw[arr] (repo) -- (provider);
\draw[arr] (provider) -- (observer);
\draw[arr] (observer) -- (state);
\draw[arr] (state) -- (ui);
\end{tikzpicture}
\end{center}

Note que a única parte deste ciclo que \emph{recompõe} algo é o último
passo — todo o caminho até ali (ViewModel, Repository, Provider,
SQLite) é Kotlin/Android comum, sem nenhuma dependência de Compose, o
que é exatamente a separação de responsabilidades que o MVVM propõe.

\section{Conclusão}

O projeto \texttt{ComposeReview} cobre, em pouco código, sete tópicos
centrais de um app Android moderno. Na Parte I, um contador e uma lista
de tarefas revisam Jetpack Compose "puro": funções \texttt{@Composable}
formam a árvore de UI; \texttt{state} (via \texttt{mutableStateOf} e
\texttt{mutableStateListOf}) guarda os dados que podem mudar;
\texttt{remember} preserva esse estado entre execuções; e a
composição/recomposição mantém a tela sincronizada com o estado, seguindo
o modelo declarativo \texttt{UI = f(state)}. Na Parte II, a tela Notas
mostra como esses mesmos conceitos de Compose se encaixam em uma
arquitetura maior: um \texttt{ContentProvider} compartilha dados
persistidos em SQLite por trás de um contrato de URIs;
\texttt{SharedPreferences} guarda preferências simples (com limitações
claras para qualquer coisa maior que isso); \texttt{Coroutines} tiram
o I/O da thread principal e transformam callbacks em \texttt{Flow};
o \texttt{ViewModel} organiza tudo em MVVM, expondo estado via
\texttt{StateFlow} (comparado a \texttt{LiveData}); e uma tela separada
com Activity/Fragment clássicos deixa visível o ciclo de vida que, em
um app só-Compose de Activity única, ficaria escondido.

\end{document}
